\documentclass[aspectratio=169,10pt,english]{beamer}

\input{../../../_preambulo_beamer.tex}
\input{../../../_comandos_beamer.tex}

\selectlanguage{english}

\title{MA1001B - Statistical Modeling for Decision Making}
\subtitle{Lesson 16: Hypergeometric Distribution: Sampling Without Replacement}
\author{}
\institute{Tecnol\'ogico de Monterrey}
\date{}

\begin{document}

\begin{frame}[plain]
  \vspace{1.2cm}
  {\Large\bfseries MA1001B - Statistical Modeling for Decision Making\par}
  \vspace{0.7cm}
  {\large Lesson 16: Hypergeometric Distribution\par}
  \vspace{0.15cm}
  {\large Sampling Without Replacement\par}
  \vspace{0.7cm}
  {\color{TecRojo}\rule{\textwidth}{1.2pt}}
  \vspace{0.8cm}

  MA1001B Faculty\\[0.3cm]
  Term\\[0.3cm]
  Tecnol\'ogico de Monterrey
\end{frame}

\begin{frame}{The Lot Is Finite}
  \begin{alertblock}{Guiding question}
    If $6$ of $80$ units are defective and $10$ are drawn without replacement,
    what is $P(X=0)$, and why is a binomial approximation too optimistic?
  \end{alertblock}

  \vspace{0.6em}
  By the end of this lesson, you can:
  \begin{itemize}
    \item identify $N$, $K$, and $n$ in a without-replacement draw;
    \item compute hypergeometric $P(X=0)$;
    \item compute $E[X]=nK/N$;
    \item explain why the binomial overstates $P(X=0)$ when $n/N$ is not small.
  \end{itemize}

  \vfill
  \scriptsize All lots used today are synthetic. No real company data are used.
\end{frame}

\begin{frame}{Hypergeometric Ingredients}
  \small
  \[
    N=80,\qquad K=6,\qquad n=10,\qquad \frac{n}{N}=0.125.
  \]
  $X$ is the number of defectives in the draw. Because units are not replaced,
  each draw changes the remaining lot, so the trials are not independent.
\end{frame}

\begin{frame}[fragile]{Step 1: $P(X=0)$ Without Replacement}
  \begin{columns}[T]
    \begin{column}{0.42\textwidth}
      \small
      \begin{lstlisting}
p0 = stats.hypergeom.pmf(
    0, M=80, n=6, N=10
)
      \end{lstlisting}

      \begin{block}{Verified output}
        $n/N=0.125000$\\
        $P(X=0)=0.436326$
      \end{block}
    \end{column}
    \begin{column}{0.58\textwidth}
      \centering
      \includegraphics[width=\textwidth]{../figures/hypergeometric_01_without_replacement.png}
      \vspace{0.2em}
      \scriptsize Hypergeometric PMF from Step 1
    \end{column}
  \end{columns}
\end{frame}

\begin{frame}{The Mean Does Not Need Independence}
  \small
  \[
    E[X]=\frac{nK}{N}=\frac{10\times 6}{80}=0.75.
  \]
  That mean matches a binomial with $p=K/N=0.075$. Matching means does not
  match $P(X=0)$.
\end{frame}

\begin{frame}[fragile]{Step 2: Hypergeometric Expectation}
  \begin{columns}[T]
    \begin{column}{0.40\textwidth}
      \small
      \begin{lstlisting}
mean = n * K / N
      \end{lstlisting}

      \begin{block}{Verified output}
        $nK/N=0.750000$\\
        scipy mean $=0.750000$
      \end{block}
    \end{column}
    \begin{column}{0.60\textwidth}
      \centering
      \includegraphics[width=\textwidth]{../figures/hypergeometric_02_expectation.png}
      \vspace{0.2em}
      \scriptsize Mean marked on the PMF from Step 2
    \end{column}
  \end{columns}
\end{frame}

\begin{frame}{Step 3: Binomial Overstates $P(X=0)$}
  \begin{columns}[T]
    \begin{column}{0.42\textwidth}
      \small
      Binomial approximation: $X\sim\mathrm{Binomial}(10,0.075)$.
      \[
        P_{\text{binom}}(X=0)=0.458582
      \]
      versus hypergeometric $0.436326$. Overstatement $0.022257$.

      \vspace{0.4em}
      \textbf{Limit:} the binomial overstates $P(X=0)$ when $n/N$ is not small.
      Lesson 17 models rare events with a Poisson rate.
    \end{column}
    \begin{column}{0.58\textwidth}
      \centering
      \includegraphics[width=\textwidth]{../figures/hypergeometric_03_binomial_approximation.png}
      \vspace{0.2em}
      \scriptsize Hypergeometric versus binomial from Step 3
    \end{column}
  \end{columns}
\end{frame}

\begin{frame}{Concept Check}
  \begin{enumerate}
    \item Why are the ten draws not independent?
    \item Compute $E[X]=nK/N$.
    \item Which model gives the larger $P(X=0)$?
    \item Why is $n/N=0.125$ a warning for the binomial approximation?
  \end{enumerate}

  \vspace{0.6em}
  \begin{block}{Connection to Lesson 17}
    The session can continue directly into the Poisson model for rare events
    in a shift.
  \end{block}

  \vfill
  \scriptsize Source: OpenStax, \textit{Introductory Business Statistics 2e},
  Section 4.1. CC BY-NC-SA.
\end{frame}

\appendix



\end{document}
